\SourceLeaf{260}%
Setzen wir diese Krümmungsmaße in die Reihe für \(\sigma\) ein, so erhalten wir den folgenden Ausdruck, der bis zu Größen sechster Ordnung (ausschließlich) genau ist:
\[
\begin{aligned}
\sigma=\tfrac{1}{2}ac\sin B\{&1+\tfrac{1}{120}\alpha(4pp-2qq+3qq'+3q'q')\\
&+\tfrac{1}{120}\beta(3pp-6qq+6qq'+3q'q')\\
&+\tfrac{1}{120}\gamma(3pp-2qq+qq'+4q'q')\}
\end{aligned}
\]
Die Genauigkeit bleibt dieselbe, wenn wir für \(p,q,q'\) die Größen \(c\sin B,c\cos B,c\cos B-a\) einsetzen. Dadurch ergibt sich
\[
\begin{aligned}
\relax[8]\quad\sigma=\tfrac{1}{2}ac\sin B\{&1+\tfrac{1}{120}\alpha(3aa+4cc-9ac\cos B)\\
&+\tfrac{1}{120}\beta(3aa+3cc-12ac\cos B)\\
&+\tfrac{1}{120}\gamma(4aa+3cc-9ac\cos B)\}
\end{aligned}
\]
Da aus dieser Gleichung alle Ausdrücke verschwunden sind, die sich auf die Linie \(AD\) beziehen, die senkrecht zu \(BC\) gezogen wurde, dürfen wir auch die Punkte \(A,B,C\) samt den zugehörigen Größen vertauschen. Damit gilt mit derselben Genauigkeit
\[
\begin{aligned}
\relax[9]\quad\sigma=\tfrac{1}{2}bc\sin A\{&1+\tfrac{1}{120}\alpha(3bb+3cc-12bc\cos A)\\
&+\tfrac{1}{120}\beta(3bb+4cc-9bc\cos A)\\
&+\tfrac{1}{120}\gamma(4bb+3cc-9bc\cos A)\}
\end{aligned}
\]
\[
\begin{aligned}
\relax[10]\quad\sigma=\tfrac{1}{2}ab\sin C\{&1+\tfrac{1}{120}\alpha(3aa+4bb-9ab\cos C)\\
&+\tfrac{1}{120}\beta(4aa+3bb-9ab\cos C)\\
&+\tfrac{1}{120}\gamma(3aa+3bb-12ab\cos C)\}
\end{aligned}
\]
\Article{26}
Es ist sehr nützlich, ein ebenes, geradliniges Dreieck zu betrachten, dessen Seiten die Längen \(a,b,c\) haben. Seine Winkel bezeichnen wir mit \(A^*,B^*,C^*\). Sie unterscheiden sich von den Winkeln des Dreiecks auf der gekrümmten Fläche, also von \(A,B,C\), um Größen zweiter Ordnung. Es lohnt sich, diese Unterschiede genau zu entwickeln. Die Rechnungen sind eher langwierig als schwierig, sodass es genügt, die wesentlichen Schritte anzugeben.

\Needspace{8\baselineskip}
Ersetzen wir in den Formeln [1], [4], [5] die auf \(B\) bezogenen Größen durch die entsprechenden Größen für \(C\), so erhalten wir Formeln für \(r'r',r'\cos\varphi',r'\sin\varphi'\). Die Entwicklung des Ausdrucks \(rr+r'r'-(q-q')^2-2r\cos\varphi\mathord{.}r'\cos\varphi'-2r\sin\varphi\mathord{.}r'\sin\varphi'\), der die Form \(=bb+cc-aa-2bc\cos A=2bc(\cos A^*-\cos A)\) annimmt, liefert zusammen mit der Entwicklung
